%% kapa: exam=1 round=22 year=2011 period=1 form=A subject=economics subject_ko=경제원론 seq=23 number=63
%% source: Q-Net 감정평가사 기출문제 (cst00309) artlSeq=5036790 file=extracted/1008/5036790/2011년 제22회 감정평가사 1차시험 기출문제/1교시.hwp
%% conversion: kordoc 4.24.51 -> tools/gen1.py (auto); review=proofread-sonnet; flags=none
\begin{question}{63}
시장실패에 관한 설명으로 옳지 않은 것은?
\begin{choices}
  \item 시장실패는 시장기능을 통하여 자원이 효율적으로 배분되지 않는 경우를 포함한다.
  \item 정부개입이 사회후생을 증대시키는데 도움을 줄 수 있다.
  \item 시장실패는 외부효과가 존재하는 경우 발생할 수 있다.
  \item 시장실패는 소유권이 명확하게 규정되지 않은 경우 발생할 수 있다.
  \item 코즈(Coase)정리에 의하면 시장실패는 시장에서 해결될 수 없다.
\end{choices}
\end{question}
\answer{5}
